\section{Artifact authority and reviewable completion}\label{sec:C23:completion}
A coding agent's final message should identify the artifact, the checks actually executed, their outcomes, and material limitations. This is not merely presentation quality. It gives the recipient the evidence needed to distinguish a proposed solution from a validated change. A confident statement that tests pass is not evidence if the trace records a skipped or failed command. The harness should generate factual execution status from authoritative records whenever possible.

Reviewable completion requires a stable artifact boundary. In version control, that can be a commit digest or patch against a known base. In a browser-created application, it can be an exported bundle and deployment identifier. In a computer-use task, it can be an authoritative application record with its permission state. The boundary varies, but it must allow the verifier to inspect the same object the agent claims to have completed.

Human review is an additional control, not an oracle. Reviewers have limited time, may anchor on the agent's explanation, and may not inspect every dependency or generated file. The agent should therefore expose high-risk changes and missing evidence rather than bury them in a long narrative. Evaluation can measure whether reviewers detect seeded defects under a fixed review budget. It should not claim safety merely because a workflow contains an approval button.

The strongest engineering design binds authority to the transition being approved. Approval of a patch does not authorize later unrelated edits; approval to prepare a deployment does not necessarily authorize publication. If the artifact changes after review, the approval's dependency digest should fail validation. A resumed session must recheck current permissions even if its history contains an earlier grant. These conditions connect coding agents to the broader authorization analysis of Chapter~\ref{ch:C29}.

The practical conclusion is that coding and computer-use competence is a property of an evidence-producing execution system. Better models expand what can be attempted. Better interfaces reduce avoidable friction. Stronger evaluators identify a larger subset of defects. Durable state and authority checks keep long tasks coherent. Each layer has an identifiable contribution and a residual failure surface. Scientific comparison requires preserving those layers rather than attributing every successful artifact to an undifferentiated agent product.
